\section{Introduction}

One of the most striking features of modern deep learning \cite{lecun2015deep} is the ability of neural networks to learn complex high-dimensional functions from data with extraordinary effectiveness \cite{sejnowski2020unreasonable}. Yet this empirical success still outpaces our theoretical understanding. 
A large body of empirical work suggests that deep networks do not merely fit input-output relations, but progressively build structured representations. In convnets, feature-visualization studies have shown that successive layers extract patterns of increasing complexity, from local edge-like detectors to more semantic motifs \cite{zeiler2014visualizing}. Transfer-learning experiments indicate that earlier layers tend to contain more general features, whereas deeper layers become increasingly specialized to the task \cite{yosinski2014how}. The {\it platonic representation} viewpoint has emphasized that models often learn surprisingly aligned representations across architectures and training procedures \cite{huh2024platonic}. These observations point to a central challenge: 

\emph{Can we understand which features are discovered during training, in which order, and at what sample complexity? And can we understand why some functions are learned more efficiently with multi-layer models?}

Several theoretical perspectives have clarified important parts of this picture. In the lazy regime \cite{chizat2019lazy}, or equivalently in the Neural Tangent Kernel (NTK) limit \cite{jacot2018neural,lee2019wide}, training is described as kernel regression in an essentially fixed feature space, giving a powerful theory of optimization and generalization but largely bypassing feature learning. Mean-field analyses of two-layer networks capture genuine parameter evolution and representation learning \cite{mei2018mean,chizat2018global,rotskoff2018neural,sirignano2020meanfield}, while tensor-program and dynamical mean-field theory approaches provide a general framework for infinite-width limits, including feature-learning regimes under suitable parametrizations \cite{yang2021tuning,hajjar2024training,bordelon2022self,pacelli2023statistical}. In parallel, stylized high-dimensional models, such as single-index and multi-index settings, have made it possible to analyze when neural networks recover latent structure from data \cite{ghorbani2020when,bietti2022learning,damian2024computational,troiani2024fundamental}, and recent work has begun to clarify the computational role of depth in synthetic toy models \cite{garnierbrun2025how,cagnetta2024how,favero2021locality,favero2025how,dandi2026computational,ren2026provable}. Still, we lack a simple predictive mechanism for how deep networks select, organize, and refine features across layers.

In this work, we propose such a mechanism, which we call {\it Neural Low-Degree Filtering} (Neural LoFi). Neural LoFi is a stylized, mathematically tractable iterative spectral surrogate for feature learning: Given a current representation, each layer forms a label-weighted moment operator on the current features, selects its leading eigendirections, and lifts the resulting projected features through a nonlinear random feature map. This procedure is motivated by a stylized small-initialization limit of gradient-based training, in which the feature-learning component of the layerwise dynamics is governed by a Hessian-like, label-weighted second-order operator. 



The analysis of Neural LoFi leads to two main lessons. First, at a fixed representation, feature learning is governed by a \emph{relevance--complexity} trade-off: the next layer selects features that have large low-degree correlation with the label while remaining simple in the geometry induced by the current representation. Neural LoFi makes this trade-off explicit through a variational problem, where relevance is measured by supervised low-degree correlation and complexity by the RKHS norm. This variational view also yields an {\it explicit, data-driven criterion for feature emergence}: a new direction becomes learnable when its population correlation rises above the empirical noise floor, whose scale is controlled by the residual effective dimension of the current kernel.


Second, depth makes this process powerful because the selected features are lifted into a new representation, changing the geometry in which the next layer searches for signal. Thus a deep network can turn structure that is high-degree in the input into structure that is low-degree in an intermediate representation. This leads to a principle of {\it low-degree compositionality}: deep learning is efficient when each stage of the target is not only compositional, but visible through low-degree correlations in the current representation. The same viewpoint also gives a kernel interpretation of feature learning: Neural LoFi constructs a sequence of task-adaptive kernels, one per layer, rather than a single fixed kernel chosen before seeing the labels. In this sense, depth allows for an adaptive multilayer kernel construction driven by supervised low-degree feature selection. 

Although we use Neural LoFi primarily as a theoretical lens, the resulting procedure is also of independent interest: it is layerwise, backpropagation-free, and based only on label-weighted low-degree moments of the current representation. In this sense, it provides a simple spectral primitive for feature discovery, sitting between fixed random-feature models, which do not adapt their representation, and full end-to-end backpropagation, whose dynamics are much harder to analyze.


The rest of the paper develops this picture. We first introduce Neural LoFi and motivate its label-weighted moment operator from the early feature-learning dynamics of gradient descent. We then derive an RKHS variational characterization, which makes explicit the relevance--complexity trade-off solved at each layer. This leads to a criterion for feature emergence, expressed through the residual effective dimension of the current kernel, and to the principle of low-degree compositionality. We conclude with a study of a solvable mathematical model that illustrates out main point (emergence of features and low-degree compositionality), and with real-data experiments, including fully connected and convolutional architectures, illustrating the predicted mechanisms arising in practice.

A public implementation, together with code to reproduce the numerical illustrations, is available at
\url{https://github.com/IdePHICS/Neural-LoFi-Theory}.
